% GENERATED FILE. Edit canonical lessons, then run npm run generate:books.
% canonical-source-hash: fnv1a64:e3f099e46b514fb1
% canonical-lessons: TA-C148-alamana, TA-C148-akalamana, TA-C148-kurukiya, TA-C148-kaynta, TA-C148-nanainta, TA-R148-first-pass-animals-and-the-body, TA-R148-second-pass-the-body-and-describing-words

\chapter{Deep, Wide, Narrow, Dry, Wet}
\label{ch:ta-deep-wide-narrow-dry-wet}
\hlchaptermodality{148}

\begin{chapteropening}
\textbf{By the end of this chapter:} \emph{I can say deep, wide, narrow, dry and wet.}

Complete the last lesson of chapter 148: i can say deep, wide, narrow, dry and wet.
\end{chapteropening}

\section[āḻamāṉa]{\ta{ஆழமான} (āḻamāṉa) — deep}

\label{lesson:TA-C148-alamana}



\begin{quote}
Before the new one: say the Tamil for lazy; a lazy person, then the Tamil for brave.
\end{quote}

\subsection*{You'll want to know: \ta{ஆழமான}}
\textbf{\ta{ஆழமான}} — \emph{āḻamāṉa} — \textquotedblleft{}deep\textquotedblright{}.

\begin{grammarlens}[title={What you've built}]
One more word for this chapter.
\end{grammarlens}

\subsection*{Your turn}
\begin{itemize}
\raggedright
  \item \emph{Say it:} \emph{āḻamāṉa}
  \item \emph{Say it:} \emph{āḻamāṉa}, once more
  \item \emph{Say it:} say \emph{āḻamāṉa}
  \item \emph{Recall:} say the Tamil for lazy; a lazy person, then the Tamil for brave, then say \emph{āḻamāṉa} again
\end{itemize}

\begin{tcolorbox}[breakable,colback=teal!4,colframe=teal!35!black,title={Before you move on}]
What does \ta{ஆழமான} mean? (\textquotedblleft{}Deep\textquotedblright{}.) Say it once more.
\end{tcolorbox}
\section[akalamāṉa]{\ta{அகலமான} (akalamāṉa) — wide}

\label{lesson:TA-C148-akalamana}



\begin{quote}
Before the new one: say the Tamil for brave, then the Tamil for deep.
\end{quote}

\subsection*{You'll want to know: \ta{அகலமான}}
\textbf{\ta{அகலமான}} — \emph{akalamāṉa} — \textquotedblleft{}wide\textquotedblright{}.

\begin{grammarlens}[title={What you've built}]
One more word for this chapter.
\end{grammarlens}

\subsection*{Your turn}
\begin{itemize}
\raggedright
  \item \emph{Say it:} \emph{akalamāṉa}
  \item \emph{Say it:} \emph{akalamāṉa}, once more
  \item \emph{Say it:} say \emph{akalamāṉa}
  \item \emph{Recall:} say the Tamil for brave, then the Tamil for deep, then say \emph{akalamāṉa} again
\end{itemize}

\begin{tcolorbox}[breakable,colback=teal!4,colframe=teal!35!black,title={Before you move on}]
What does \ta{அகலமான} mean? (\textquotedblleft{}Wide\textquotedblright{}.) Say it once more.
\end{tcolorbox}
\section[kuṟukiya]{\ta{குறுகிய} (kuṟukiya) — narrow}

\label{lesson:TA-C148-kurukiya}



\begin{quote}
Before the new one: say the Tamil for deep, then the Tamil for wide.
\end{quote}

\subsection*{You'll want to know: \ta{குறுகிய}}
\textbf{\ta{குறுகிய}} — \emph{kuṟukiya} — \textquotedblleft{}narrow\textquotedblright{}.

\begin{grammarlens}[title={What you've built}]
One more word for this chapter.
\end{grammarlens}

\subsection*{Your turn}
\begin{itemize}
\raggedright
  \item \emph{Say it:} \emph{kuṟukiya}
  \item \emph{Say it:} \emph{kuṟukiya}, once more
  \item \emph{Say it:} say \emph{kuṟukiya}
  \item \emph{Recall:} say the Tamil for deep, then the Tamil for wide, then say \emph{kuṟukiya} again
\end{itemize}

\begin{tcolorbox}[breakable,colback=teal!4,colframe=teal!35!black,title={Before you move on}]
What does \ta{குறுகிய} mean? (\textquotedblleft{}Narrow\textquotedblright{}.) Say it once more.
\end{tcolorbox}
\section[kāynta]{\ta{காய்ந்த} (kāynta) — dry}

\label{lesson:TA-C148-kaynta}



\begin{quote}
Before the new one: say the Tamil for wide, then the Tamil for narrow.
\end{quote}

\subsection*{You'll want to know: \ta{காய்ந்த}}
\textbf{\ta{காய்ந்த}} — \emph{kāynta} — \textquotedblleft{}dry\textquotedblright{}.

\begin{grammarlens}[title={What you've built}]
One more word for this chapter.
\end{grammarlens}

\subsection*{Your turn}
\begin{itemize}
\raggedright
  \item \emph{Say it:} \emph{kāynta}
  \item \emph{Say it:} \emph{kāynta}, once more
  \item \emph{Say it:} say \emph{kāynta}
  \item \emph{Recall:} say the Tamil for wide, then the Tamil for narrow, then say \emph{kāynta} again
\end{itemize}

\begin{tcolorbox}[breakable,colback=teal!4,colframe=teal!35!black,title={Before you move on}]
What does \ta{காய்ந்த} mean? (\textquotedblleft{}Dry\textquotedblright{}.) Say it once more.
\end{tcolorbox}
\section[naṉainta]{\ta{நனைந்த} (naṉainta) — wet}

\label{lesson:TA-C148-nanainta}



\begin{quote}
Before the new one: say the Tamil for narrow, then the Tamil for dry.
\end{quote}

\subsection*{You'll want to know: \ta{நனைந்த}}
\textbf{\ta{நனைந்த}} — \emph{naṉainta} — \textquotedblleft{}wet\textquotedblright{}.

\begin{grammarlens}[title={What you've built}]
One more word for this chapter.
\end{grammarlens}

\subsection*{Your turn}
\begin{itemize}
\raggedright
  \item \emph{Say it:} \emph{naṉainta}
  \item \emph{Say it:} \emph{naṉainta}, once more
  \item \emph{Say it:} say \emph{naṉainta}
  \item \emph{Recall:} say the Tamil for narrow, then the Tamil for dry, then say \emph{naṉainta} again
\end{itemize}

\begin{tcolorbox}[breakable,colback=teal!4,colframe=teal!35!black,title={Before you move on}]
What does \ta{நனைந்த} mean? (\textquotedblleft{}Wet\textquotedblright{}.) Say it once more.
\end{tcolorbox}
\section[(dialogue)]{Animals and the body — a first pass}

\label{lesson:TA-R148-first-pass-animals-and-the-body}



\begin{quote}
Say the words for dry and wet.
\end{quote}

\subsection*{Your turn}
Say these aloud:

\begin{itemize}
\raggedright
  \item grapes, a pomegranate, a pickle, a waterfall, a hill
  \item an island, an insect, a spider, a pig, a donkey
  \item a camel, a duck, a cheek, a wrist, an elbow
  \item a thumb, bad, beautiful, long, short
\end{itemize}

\begin{tcolorbox}[breakable,colback=teal!4,colframe=teal!35!black,title={Before you move on}]
Grapes? (\textbf{\ta{திராட்சை}}, \emph{tirāṭcai}.) A pomegranate? (\textbf{\ta{மாதுளை}}, \emph{mātuḷai}.) A pickle? (\textbf{\ta{ஊறுகாய்}}, \emph{ūṟukāy}.) A waterfall? (\textbf{\ta{அருவி}}, \emph{aruvi}.) A hill? (\textbf{\ta{குன்று}}, \emph{kuṉṟu}.) An island? (\textbf{\ta{தீவு}}, \emph{tīvu}.) An insect? (\textbf{\ta{பூச்சி}}, \emph{pūcci}.) A spider? (\textbf{\ta{சிலந்தி}}, \emph{cilanti}.) A pig? (\textbf{\ta{பன்றி}}, \emph{paṉṟi}.) A donkey? (\textbf{\ta{கழுதை}}, \emph{kaḻutai}.) A camel? (\textbf{\ta{ஒட்டகம்}}, \emph{ottakam}.) A duck? (\textbf{\ta{வாத்து}}, \emph{vāttu}.) A cheek? (\textbf{\ta{கன்னம்}}, \emph{kaṉṉam}.) A wrist? (\textbf{\ta{மணிக்கட்டு}}, \emph{maṇikkaṭṭu}.) An elbow? (\textbf{\ta{முழங்கை}}, \emph{muḻaṅkai}.) A thumb? (\textbf{\ta{கட்டைவிரல்}}, \emph{kaṭṭaiviral}.) Bad? (\textbf{\ta{கெட்ட}}, \emph{keṭṭa}.) Beautiful? (\textbf{\ta{அழகான}}, \emph{aḻakāṉa}.) Long? (\textbf{\ta{நீளமான}}, \emph{nīḷamāṉa}.) Short? (\textbf{\ta{குட்டையான}}, \emph{kuṭṭaiyāṉa}.)
\end{tcolorbox}
\section[(dialogue)]{The body and describing words — a second pass}

\label{lesson:TA-R148-second-pass-the-body-and-describing-words}



\begin{quote}
Say the words for tall and hot.
\end{quote}

\subsection*{Your turn}
Say these aloud:

\begin{itemize}
\raggedright
  \item tall, hot, cold, heavy, light
  \item soft, hard, clean, dirty, fast
  \item bitter, sour, spicy, rich, poor
  \item young, old, clever; a clever person, lazy; a lazy person, brave
  \item deep, wide, narrow, dry, wet
\end{itemize}

\begin{tcolorbox}[breakable,colback=teal!4,colframe=teal!35!black,title={Before you move on}]
Tall? (\textbf{\ta{உயரமான}}, \emph{uyaramāṉa}.) Hot? (\textbf{\ta{சூடான}}, \emph{cūṭāṉa}.) Cold? (\textbf{\ta{குளிர்ந்த}}, \emph{kuḷirnta}.) Heavy? (\textbf{\ta{கனமான}}, \emph{kaṉamāṉa}.) Light? (\textbf{\ta{லேசான}}, \emph{lēcāṉa}.) Soft? (\textbf{\ta{மென்மையான}}, \emph{meṉmaiyāṉa}.) Hard? (\textbf{\ta{கடினமான}}, \emph{kaṭiṉamāṉa}.) Clean? (\textbf{\ta{சுத்தமான}}, \emph{cuttamāṉa}.) Dirty? (\textbf{\ta{அழுக்கான}}, \emph{aḻukkāṉa}.) Fast? (\textbf{\ta{வேகமான}}, \emph{vēkamāṉa}.) Bitter? (\textbf{\ta{கசப்பு}}, \emph{kacappu}.) Sour? (\textbf{\ta{புளிப்பு}}, \emph{puḷippu}.) Spicy? (\textbf{\ta{காரம்}}, \emph{kāram}.) Rich? (\textbf{\ta{பணக்கார}}, \emph{paṇakkāra}.) Poor? (\textbf{\ta{ஏழை}}, \emph{ēḻai}.) Young? (\textbf{\ta{இளைய}}, \emph{iḷaiya}.) Old? (\textbf{\ta{வயதான}}, \emph{vayatāṉa}.) Clever; a clever person? (\textbf{\ta{புத்திசாலி}}, \emph{putticāli}.) Lazy; a lazy person? (\textbf{\ta{சோம்பேறி}}, \emph{cōmpēṟi}.) Brave? (\textbf{\ta{தைரியமான}}, \emph{tairiyamāṉa}.) Deep? (\textbf{\ta{ஆழமான}}, \emph{āḻamāṉa}.) Wide? (\textbf{\ta{அகலமான}}, \emph{akalamāṉa}.) Narrow? (\textbf{\ta{குறுகிய}}, \emph{kuṟukiya}.) Dry? (\textbf{\ta{காய்ந்த}}, \emph{kāynta}.) Wet? (\textbf{\ta{நனைந்த}}, \emph{naṉainta}.)
\end{tcolorbox}
